%! Setting Up a Test for the Difference of Two Population Proportions — Unit 6, Lesson 6.10 — Group Activity
\documentclass[10pt]{article}
\usepackage{apstats-article}
\usepackage{apstats-boxes}
\newcommand{\TallMath}[1]{$\displaystyle #1\rule[-1.4em]{0pt}{3.2em}$}

\begin{document}

\pageheader{Unit 6, Lesson 6.10}{Group Activity}
\namepartnerperiod

\vspace{0.15in}

\textbf{Directions:} For each scenario, complete the \textit{State} (hypotheses) and
\textit{Plan} (conditions) steps only.

\begin{scenariobox}[Scenario 1 --- Online vs.\ In-Person Test Preparation]{sky}
A researcher wants to test whether the proportion of students who pass a licensing exam
differs between those who prepared online and those who prepared in person.
Online: 84 of 140 passed. In-person: 105 of 160 passed.

\begin{enumerate}[label=\arabic*.]
  \item Define $p_1$ and $p_2$ in context. \writeline

  \item State $H_0$ and $H_a$. Is this one-sided or two-sided? \writelines{2}

  \item Compute $\hat p_1$, $\hat p_2$, and $\hat p_c$. \writelines{2}

  \item Verify all three conditions using $\hat p_c$ for Large Counts.
        \writelines{4}
\end{enumerate}
\end{scenariobox}

\begin{scenariobox}[Scenario 2 --- Seatbelt Compliance Among Teen Drivers]{sky}
A traffic safety officer hypothesizes that a higher proportion of 16-year-olds than
17-year-olds fail to wear a seat belt. In a random roadside check:
16-year-olds: 18 of 90 not wearing a seat belt.
17-year-olds: 12 of 110 not wearing a seat belt.

\begin{enumerate}[label=\arabic*.]
  \item Define $p_{16}$ and $p_{17}$ and state $H_0$ and $H_a$ (one-sided). \writelines{2}

  \item Compute $\hat p_{16}$, $\hat p_{17}$, and $\hat p_c$. \writelines{2}

  \item Verify conditions. \writelines{3}
\end{enumerate}
\end{scenariobox}

\end{document}
